
\begin{frame}[fragile]{L'opérateur qui nous manque}

La jointure: opération très courante, potentiellement coûteuse.

\vfill

L'opérateur de jointure complète notre petit catalogue pour (presque) toutes 
les requêtes SQL.

\vfill

\begin{minted}[frame=leftline,
               baselinestretch=.9,
               fontsize=\small,
               framesep=2mm]{sql}
    select a1, a2, .., an
    from T1, T2, ..., Tm
    where T1.x = T2.y and ... 
    order by ...
\end{minted}

\vfill

Plusieurs algorithmes possibles.


\end{frame}

\begin{frame}{Principaux algorithmes}

On ve se limiter à quelques exemples représentatifs:

\vfill


\alert{Jointure avec index}

\begin{redbullet}

  \item Algorithme de jointure par boucles imbriquées indexées.

\end{redbullet}

\vfill

\alert{Jointure sans index}

\vfill
\begin{blocgauche}
\begin{redbullet}

  \item Le plus simple\,: \alert{boucles imbriquées (non indexée)}.

   \item Plus sophistiqué\,: la \alert{jointure par hachage}.
\end{redbullet}
\end{blocgauche}
\end{frame}

%%%%%%%%%%%%%%%%%%%%%%%%%%%%%%%%%%%%%%%%%%%%%%%%%%%%%%%
\begin{frame}[fragile]{Jointure avec index}

Très \alert{courant}\,; on effectue
\alert{naturellement} la jointure sur les clés primaires/étrangères.

\begin{redbullet}

 \item Les films et leur metteur en scène
 
\begin{minted}[frame=leftline,
               baselinestretch=.9,
               fontsize=\footnotesize,
               framesep=2mm]{sql}
      select * from Film as f, Artiste as a 
      where f.id_realisateur = a.id
\end{minted}

      
 \item Les employés et leur département

\begin{minted}[frame=leftline,
               baselinestretch=.9,
               fontsize=\footnotesize,
               framesep=2mm]{sql}
      select * from emp e, dept d
      where e.dnum = d.num
\end{minted}

\end{redbullet}

Garantit \alert{au moins} un index sur la condition de jointure.
\end{frame}



\begin{frame}{Jointure avec index: l'algorithme}


\only<1-1|handout:1>{
On parcourt séquentiellement la table contenant la clé étrangère.


\vfill
\figSlideWithSize{planEx-indexedjoin-1}{5cm}
\vfill
On obtient des nuplets \texttt{employé}, avec leur no de département.
}
\only<2-2|handout:2>{
On utilise le no de département pour un accès à l'index \texttt{Dept}.

\vfill
\figSlideWithSize{planEx-indexedjoin-2}{7cm}
\vfill
On obtient des paires \texttt{[employé, adrDept]}.
}
\only<3-3|handout:3>{
Il reste à résoudre l'adresse du département avec un accès direct.

\vfill
\figSlideWithSize{planEx-indexedjoin-3}{7cm}
\vfill
On obtient des paires \texttt{[employé, dept]}.
}

\only<4-4|handout:4>{
Avantages\,:
  \begin{redbullet}
      \item  Efficace (un parcours, plus des recherches par adresse)

      \item Favorise le temps de réponse \textbf{et} le temps
                 d'exécution
\end{redbullet}
}

\end{frame}

%%%%%%%%%%%%%%%%%%%%%%%%%%%%%%%%%%%%%%%%%%%%%%%%%%%%%%%

\begin{frame}[fragile]{Jointures par boucles imbriquées}

Pas d'index\,? La méthode de base est d'énumérer \alert{toutes}
les solutions possibles.

\vfill

\figSlideWithSize{joinlist}{10cm}


\vfill

\alert{Coût quadratique}. 
Acceptable pour
deux petites tables.
\end{frame}



\begin{frame}[fragile]{Jointures par hachage}

Meilleure solution: construire une table de hachage sur une des tables.

\vfill

\figSlideWithSize{hashjoinlist}{9cm}


\vfill

Evite les $O(n^2)$ comparaisons. Appelons cette méthode \alert{JoinList}.
\end{frame}

\begin{frame}{Mémoire insuffisante?}

Essayons de placer \alert{une} des deux tables en mémoire.

\vfill


\figSlideWithSize{nestedloop-improved}{8cm}

\vfill

\begin{tabular}{p{9cm}p{5cm}}

On charge l'autre bloc par bloc; on applique  \alert{JoinList}.

\vfill

Une seule lecture de chaque table suffit.
&
~
\end{tabular}

.

\end{frame}


\begin{frame}{Et quand \emph{aucune} table ne tient en mémoire?}


\only<1-1|handout:1>{

 On hache la plus petite des deux tables en $k$ fragments.
\vfill

\figSlideWithSize{hashjoin-1}{9cm}

\vfill

\alert{Essentiel}\,: les fragments
doivent tenir, chacun, en mémoire.
}


\only<2-2|handout:2-2>{

 On hache la seconde table, avec la même fonction $h()$, 
             en $k$ autres fragments.
\vfill

\figSlideWithSize{hashjoin}{9cm}

\vfill

\begin{tabular}{p{9cm}p{5cm}}

Cette fois, on n'impose pas la contrainte que les fragments tiennent en mémoire.
&
~
\end{tabular}
}
\only<3-3|handout:3-3>{
\alert{On effectue
la jointure sur les paires de fragments $(F_1^R, F_1^S), (F_2^R, F_2^S),
 (F_k^R, F_k^S),$}

\vfill

\figSlideWithSize{hashjoin}{9cm}

\vfill
\begin{tabular}{p{9cm}p{5cm}}
\alert{Propriété}: Deux nuplets  $r$ et  $s$
doivent être joints si et seulement si ils sont dans des fragments associés.
&
~
\end{tabular}

}
\end{frame}



%%%%%%%%%%%%%%%%%%%%%%%%%%%%%%%%%%%%%%%%%%%%%%%%%%%%%%%
\begin{frame}{Illustration\,: phase de jointure}

On charge  $F_R^i$ de $R$ en mémoire; on parcourt 
$F_S^i$ de $S$ et on joint.

\vfill

\figSlideWithSize{hashjoin2}{9cm}

\vfill
\begin{tabular}{p{9cm}p{5cm}}
\alert{Déjà vu?} Oui: jointure par boucles imbriquées
quand une table tient en mémoire.
&
~
\end{tabular}

\end{frame}

%%%%%%%%%%%%%%%%%%%%%%%%%%%%%%%%%%%%%%%%%%%%%%%%%%%%%%%
\begin{frame}{Résumé: la jointure}

Un opérateur potentiellement coûteux. Quelques principes généraux:

\vfill

\begin{redbullet}

   \item Si \alert{une} table tient en mémoire: jointure par boucle imbriquées, ou hachage.

\item
 Si \alert{au moins un} index est utilisable: jointure par boucle imbriquées indexée.

\item
Si une des deux tables beaucoup plus petite que l'autre: jointure par hachage.

\item

Sinon: jointure par tri-fusion (non présenté).

\end{redbullet}

\vfill

\begin{tabular}{p{9cm}p{5cm}}

\alert{Décision très complexe}, prise par la système en fonction des statistiques.

&
~
\end{tabular}


\only<2|handout:2>{
\centerline{\alert{Merci!}}
}

\end{frame}